% !TEX root = ../main.tex

\section{Literature Review}

Industrial safety management has increasingly relied on artificial intelligence to reduce the frequency of operator-related accidents on heavy machinery and construction sites~\cite{ref1}. Fatigue detection has received considerable attention in this area, where facial landmark analysis is used to derive indicators such as eye aspect ratio (EAR), percentage of eyelid closure over time (PERCLOS), and head pose, which are subsequently classified using machine learning models such as XGBoost or support vector machines~\cite{ref2}. These approaches are generally effective at identifying drowsiness from behavioral cues, although their reliability depends heavily on lighting conditions and camera placement inside the cabin.

Alongside fatigue monitoring, vision-based object detection has been applied for blind-spot and worker-proximity awareness around excavators, loaders, and similar equipment. Architectures from the YOLO family are commonly used for this purpose, with detected objects mapped onto predefined danger zones around the machine to trigger proximity alerts~\cite{ref3}. Separately, a body of work has used historical incident data, including OSHA-reported injury records, to train predictive models for contextual risk estimation. These models typically use gradient-boosted classifiers to relate variables such as machine type, task category, weather, and time of day to accident likelihood, producing a risk score that supports supervisory decision-making rather than real-time intervention~\cite{ref4}.

Reinforcement learning has also been explored for adaptive decision-making in safety-related systems, most notably through policy-gradient methods such as Proximal Policy Optimization (PPO), trained within simulated environments to learn intervention behavior rather than relying on manually specified rules~\cite{ref5}. However, the majority of this work has been directed at autonomous vehicle or robotic control problems, with limited application to systems that must warn a human operator rather than control the machine directly. In parallel, explainability techniques such as SHAP have gained traction in safety-critical domains, where they are used to attribute model outputs to the underlying input features and support operator or supervisor trust in automated alerts~\cite{ref6}.

Taken together, these lines of work show that fatigue detection, proximity-based hazard detection, contextual risk estimation, and reinforcement-learning-based decision-making have each reached a reasonable level of maturity when studied independently. What is comparatively underexplored is their combination into a single system capable of reasoning jointly about operator condition, external hazards, and machine state, particularly one that also accounts for how the resulting alert should be delivered to a specific operator.

\subsection{Limitations and Research Gap}

While the reviewed techniques have individually contributed to improved hazard detection and risk estimation, several limitations remain when these systems are considered as a whole. Most existing systems are designed around a single hazard category, so a given deployment typically handles either fatigue monitoring or proximity detection, but rarely both together with additional factors such as machine tilt or stability. Where intervention logic is present, it is usually governed by fixed thresholds that do not adjust to individual operator behavior or to the operating context, which can result in either delayed warnings or an excess of false alerts. Although reinforcement learning has been used for adaptive decision-making in related domains, most reported implementations do not include an explicit mechanism to guarantee that a learned policy cannot under-respond to a genuinely critical situation, since safety behavior is implicit in the reward formulation rather than structurally enforced. Explainability, similarly, is usually applied during offline model evaluation rather than surfaced to the end user at the moment an alert is triggered, which limits an operator's ability to trust or verify the decision in real time.

A further limitation concerns personalization of alert delivery. Existing systems generally assume a uniform operator profile, with alert modality fixed at design time regardless of hearing or visual impairment, experience level, or language, despite this directly affecting whether a warning is perceived and understood in time. Where multiple AI modules are deployed together, they typically operate as separate pipelines rather than a coordinated decision system, leaving the task of synthesizing several alert streams to a human supervisor. Finally, a substantial portion of the reviewed literature assumes network or cloud connectivity for inference, which is not a realistic assumption for remote worksites and heavy-machinery cabs.

These gaps indicate a clear need for an integrated, priority-aware safety framework capable of combining multi-modal hazard sensing, a hybrid rule-based and learning-based decision architecture, accessibility-driven personalization, and real-time explainability, while operating fully offline. To address these gaps, this work proposes SAARTHI, a unified safety copilot that fuses fatigue detection, blind-spot monitoring, and machine tilt sensing with a two-tier decision engine --- deterministic hard rules for known critical hazards, combined with a PPO-trained adaptive policy for softer intervention levels --- together with accessibility-aware alert personalization and reason-based explainability, deployed entirely on-device.
